%% template.tex
%% from
%% bare_conf.tex
%% V1.4b
%% 2015/08/26
%% by Michael Shell
%% See:
%% http://www.michaelshell.org/
%% for current contact information.
%%
%% This is a skeleton file demonstrating the use of IEEEtran.cls
%% (requires IEEEtran.cls version 1.8b or later) with an IEEE
%% conference paper.
%%
%% Support sites:
%% http://www.michaelshell.org/tex/ieeetran/
%% http://www.ctan.org/pkg/ieeetran
%% and
%% http://www.ieee.org/

%%*************************************************************************
%% Legal Notice:
%% This code is offered as-is without any warranty either expressed or
%% implied; without even the implied warranty of MERCHANTABILITY or
%% FITNESS FOR A PARTICULAR PURPOSE!
%% User assumes all risk.
%% In no event shall the IEEE or any contributor to this code be liable for
%% any damages or losses, including, but not limited to, incidental,
%% consequential, or any other damages, resulting from the use or misuse
%% of any information contained here.
%%
%% All comments are the opinions of their respective authors and are not
%% necessarily endorsed by the IEEE.
%%
%% This work is distributed under the LaTeX Project Public License (LPPL)
%% ( http://www.latex-project.org/ ) version 1.3, and may be freely used,
%% distributed and modified. A copy of the LPPL, version 1.3, is included
%% in the base LaTeX documentation of all distributions of LaTeX released
%% 2003/12/01 or later.
%% Retain all contribution notices and credits.
%% ** Modified files should be clearly indicated as such, including  **
%% ** renaming them and changing author support contact information. **
%%*************************************************************************


% *** Authors should verify (and, if needed, correct) their LaTeX system  ***
% *** with the testflow diagnostic prior to trusting their LaTeX platform ***
% *** with production work. The IEEE's font choices and paper sizes can   ***
% *** trigger bugs that do not appear when using other class files.       ***                          ***
% The testflow support page is at:
% http://www.michaelshell.org/tex/testflow/

\documentclass[conference,final,]{IEEEtran}
% Some Computer Society conferences also require the compsoc mode option,
% but others use the standard conference format.
%
% If IEEEtran.cls has not been installed into the LaTeX system files,
% manually specify the path to it like:
% \documentclass[conference]{../sty/IEEEtran}





% Some very useful LaTeX packages include:
% (uncomment the ones you want to load)


% *** MISC UTILITY PACKAGES ***
%
%\usepackage{ifpdf}
% Heiko Oberdiek's ifpdf.sty is very useful if you need conditional
% compilation based on whether the output is pdf or dvi.
% usage:
% \ifpdf
%   % pdf code
% \else
%   % dvi code
% \fi
% The latest version of ifpdf.sty can be obtained from:
% http://www.ctan.org/pkg/ifpdf
% Also, note that IEEEtran.cls V1.7 and later provides a builtin
% \ifCLASSINFOpdf conditional that works the same way.
% When switching from latex to pdflatex and vice-versa, the compiler may
% have to be run twice to clear warning/error messages.






% *** CITATION PACKAGES ***
%
%\usepackage{cite}
% cite.sty was written by Donald Arseneau
% V1.6 and later of IEEEtran pre-defines the format of the cite.sty package
% \cite{} output to follow that of the IEEE. Loading the cite package will
% result in citation numbers being automatically sorted and properly
% "compressed/ranged". e.g., [1], [9], [2], [7], [5], [6] without using
% cite.sty will become [1], [2], [5]--[7], [9] using cite.sty. cite.sty's
% \cite will automatically add leading space, if needed. Use cite.sty's
% noadjust option (cite.sty V3.8 and later) if you want to turn this off
% such as if a citation ever needs to be enclosed in parenthesis.
% cite.sty is already installed on most LaTeX systems. Be sure and use
% version 5.0 (2009-03-20) and later if using hyperref.sty.
% The latest version can be obtained at:
% http://www.ctan.org/pkg/cite
% The documentation is contained in the cite.sty file itself.






% *** GRAPHICS RELATED PACKAGES ***
%
\ifCLASSINFOpdf
  % \usepackage[pdftex]{graphicx}
  % declare the path(s) where your graphic files are
  % \graphicspath{{../pdf/}{../jpeg/}}
  % and their extensions so you won't have to specify these with
  % every instance of \includegraphics
  % \DeclareGraphicsExtensions{.pdf,.jpeg,.png}
\else
  % or other class option (dvipsone, dvipdf, if not using dvips). graphicx
  % will default to the driver specified in the system graphics.cfg if no
  % driver is specified.
  % \usepackage[dvips]{graphicx}
  % declare the path(s) where your graphic files are
  % \graphicspath{{../eps/}}
  % and their extensions so you won't have to specify these with
  % every instance of \includegraphics
  % \DeclareGraphicsExtensions{.eps}
\fi
% graphicx was written by David Carlisle and Sebastian Rahtz. It is
% required if you want graphics, photos, etc. graphicx.sty is already
% installed on most LaTeX systems. The latest version and documentation
% can be obtained at:
% http://www.ctan.org/pkg/graphicx
% Another good source of documentation is "Using Imported Graphics in
% LaTeX2e" by Keith Reckdahl which can be found at:
% http://www.ctan.org/pkg/epslatex
%
% latex, and pdflatex in dvi mode, support graphics in encapsulated
% postscript (.eps) format. pdflatex in pdf mode supports graphics
% in .pdf, .jpeg, .png and .mps (metapost) formats. Users should ensure
% that all non-photo figures use a vector format (.eps, .pdf, .mps) and
% not a bitmapped formats (.jpeg, .png). The IEEE frowns on bitmapped formats
% which can result in "jaggedy"/blurry rendering of lines and letters as
% well as large increases in file sizes.
%
% You can find documentation about the pdfTeX application at:
% http://www.tug.org/applications/pdftex

\usepackage{graphicx}

% *** MATH PACKAGES ***
%
%\usepackage{amsmath}
% A popular package from the American Mathematical Society that provides
% many useful and powerful commands for dealing with mathematics.
%
% Note that the amsmath package sets \interdisplaylinepenalty to 10000
% thus preventing page breaks from occurring within multiline equations. Use:
%\interdisplaylinepenalty=2500
% after loading amsmath to restore such page breaks as IEEEtran.cls normally
% does. amsmath.sty is already installed on most LaTeX systems. The latest
% version and documentation can be obtained at:
% http://www.ctan.org/pkg/amsmath





% *** SPECIALIZED LIST PACKAGES ***
%
%\usepackage{algorithmic}
% algorithmic.sty was written by Peter Williams and Rogerio Brito.
% This package provides an algorithmic environment fo describing algorithms.
% You can use the algorithmic environment in-text or within a figure
% environment to provide for a floating algorithm. Do NOT use the algorithm
% floating environment provided by algorithm.sty (by the same authors) or
% algorithm2e.sty (by Christophe Fiorio) as the IEEE does not use dedicated
% algorithm float types and packages that provide these will not provide
% correct IEEE style captions. The latest version and documentation of
% algorithmic.sty can be obtained at:
% http://www.ctan.org/pkg/algorithms
% Also of interest may be the (relatively newer and more customizable)
% algorithmicx.sty package by Szasz Janos:
% http://www.ctan.org/pkg/algorithmicx




% *** ALIGNMENT PACKAGES ***
%
%\usepackage{array}
% Frank Mittelbach's and David Carlisle's array.sty patches and improves
% the standard LaTeX2e array and tabular environments to provide better
% appearance and additional user controls. As the default LaTeX2e table
% generation code is lacking to the point of almost being broken with
% respect to the quality of the end results, all users are strongly
% advised to use an enhanced (at the very least that provided by array.sty)
% set of table tools. array.sty is already installed on most systems. The
% latest version and documentation can be obtained at:
% http://www.ctan.org/pkg/array


% IEEEtran contains the IEEEeqnarray family of commands that can be used to
% generate multiline equations as well as matrices, tables, etc., of high
% quality.




% *** SUBFIGURE PACKAGES ***
%\ifCLASSOPTIONcompsoc
%  \usepackage[caption=false,font=normalsize,labelfont=sf,textfont=sf]{subfig}
%\else
%  \usepackage[caption=false,font=footnotesize]{subfig}
%\fi
% subfig.sty, written by Steven Douglas Cochran, is the modern replacement
% for subfigure.sty, the latter of which is no longer maintained and is
% incompatible with some LaTeX packages including fixltx2e. However,
% subfig.sty requires and automatically loads Axel Sommerfeldt's caption.sty
% which will override IEEEtran.cls' handling of captions and this will result
% in non-IEEE style figure/table captions. To prevent this problem, be sure
% and invoke subfig.sty's "caption=false" package option (available since
% subfig.sty version 1.3, 2005/06/28) as this is will preserve IEEEtran.cls
% handling of captions.
% Note that the Computer Society format requires a larger sans serif font
% than the serif footnote size font used in traditional IEEE formatting
% and thus the need to invoke different subfig.sty package options depending
% on whether compsoc mode has been enabled.
%
% The latest version and documentation of subfig.sty can be obtained at:
% http://www.ctan.org/pkg/subfig




% *** FLOAT PACKAGES ***
%
\usepackage{dblfloatfix}

%\usepackage{fixltx2e}
% fixltx2e, the successor to the earlier fix2col.sty, was written by
% Frank Mittelbach and David Carlisle. This package corrects a few problems
% in the LaTeX2e kernel, the most notable of which is that in current
% LaTeX2e releases, the ordering of single and double column floats is not
% guaranteed to be preserved. Thus, an unpatched LaTeX2e can allow a
% single column figure to be placed prior to an earlier double column
% figure.
% Be aware that LaTeX2e kernels dated 2015 and later have fixltx2e.sty's
% corrections already built into the system in which case a warning will
% be issued if an attempt is made to load fixltx2e.sty as it is no longer
% needed.
% The latest version and documentation can be found at:
% http://www.ctan.org/pkg/fixltx2e


%\usepackage{stfloats}
% stfloats.sty was written by Sigitas Tolusis. This package gives LaTeX2e
% the ability to do double column floats at the bottom of the page as well
% as the top. (e.g., "\begin{figure*}[!b]" is not normally possible in
% LaTeX2e). It also provides a command:
%\fnbelowfloat
% to enable the placement of footnotes below bottom floats (the standard
% LaTeX2e kernel puts them above bottom floats). This is an invasive package
% which rewrites many portions of the LaTeX2e float routines. It may not work
% with other packages that modify the LaTeX2e float routines. The latest
% version and documentation can be obtained at:
% http://www.ctan.org/pkg/stfloats
% Do not use the stfloats baselinefloat ability as the IEEE does not allow
% \baselineskip to stretch. Authors submitting work to the IEEE should note
% that the IEEE rarely uses double column equations and that authors should try
% to avoid such use. Do not be tempted to use the cuted.sty or midfloat.sty
% packages (also by Sigitas Tolusis) as the IEEE does not format its papers in
% such ways.
% Do not attempt to use stfloats with fixltx2e as they are incompatible.
% Instead, use Morten Hogholm'a dblfloatfix which combines the features
% of both fixltx2e and stfloats:
%
% \usepackage{dblfloatfix}
% The latest version can be found at:
% http://www.ctan.org/pkg/dblfloatfix




% *** PDF, URL AND HYPERLINK PACKAGES ***
%
%\usepackage{url}
% url.sty was written by Donald Arseneau. It provides better support for
% handling and breaking URLs. url.sty is already installed on most LaTeX
% systems. The latest version and documentation can be obtained at:
% http://www.ctan.org/pkg/url
% Basically, \url{my_url_here}.




% *** Do not adjust lengths that control margins, column widths, etc. ***
% *** Do not use packages that alter fonts (such as pslatex).         ***
% There should be no need to do such things with IEEEtran.cls V1.6 and later.
% (Unless specifically asked to do so by the journal or conference you plan
% to submit to, of course. )



%% BEGIN MY ADDITIONS %%

\usepackage[numbers]{natbib}
\bibliographystyle{IEEEtran}

\usepackage[unicode=true]{hyperref}

\hypersetup{
            pdftitle={Exploratory Visualization of Fraud Risk Patterns in Bank Transaction Data},
            pdfkeywords={bank transaction fraud, exploratory data
analysis, data visualization, class imbalance, fraud rate},
            pdfborder={0 0 0},
            breaklinks=true}
\urlstyle{same}  % don't use monospace font for urls

% Pandoc toggle for numbering sections (defaults to be off)
\setcounter{secnumdepth}{0}


% tightlist command for lists without linebreak
\providecommand{\tightlist}{%
  \setlength{\itemsep}{0pt}\setlength{\parskip}{0pt}}



\usepackage[section]{placeins}
\usepackage{booktabs}
\usepackage{booktabs}
\usepackage{longtable}
\usepackage{array}
\usepackage{multirow}
\usepackage{wrapfig}
\usepackage{float}
\usepackage{colortbl}
\usepackage{pdflscape}
\usepackage{tabu}
\usepackage{threeparttable}
\usepackage{threeparttablex}
\usepackage[normalem]{ulem}
\usepackage{makecell}
\usepackage{xcolor}

%% END MY ADDITIONS %%


\hyphenation{op-tical net-works semi-conduc-tor}

\begin{document}
%
% paper title
% Titles are generally capitalized except for words such as a, an, and, as,
% at, but, by, for, in, nor, of, on, or, the, to and up, which are usually
% not capitalized unless they are the first or last word of the title.
% Linebreaks \\ can be used within to get better formatting as desired.
% Do not put math or special symbols in the title.
\title{Exploratory Visualization of Fraud Risk Patterns in Bank
Transaction Data}

% author names and affiliations
% use a multiple column layout for up to three different
% affiliations

\author{

%% ---- classic IEEETrans wide authors' list ----------------
\IEEEauthorblockN{
%% 
Ngoc Khanh Vy Huynh, Student Number: 0908036
}
\IEEEauthorblockN{
Keyao Li, Student Number: 0895585
%%
}


%% ----------------------------------------------------------

%% ---- classic IEEETrans one column per institution --------
 %%
%% ----------------------------------------------------------





%% ---- one column per author, classic/default IEEETrans ----

%% ----------------------------------------------------------

}

% conference papers do not typically use \thanks and this command
% is locked out in conference mode. If really needed, such as for
% the acknowledgment of grants, issue a \IEEEoverridecommandlockouts
% after \documentclass

% for over three affiliations, or if they all won't fit within the width
% of the page, use this alternative format:
%
%\author{\IEEEauthorblockN{Michael Shell\IEEEauthorrefmark{1},
%Homer Simpson\IEEEauthorrefmark{2},
%James Kirk\IEEEauthorrefmark{3},
%Montgomery Scott\IEEEauthorrefmark{3} and
%Eldon Tyrell\IEEEauthorrefmark{4}}
%\IEEEauthorblockA{\IEEEauthorrefmark{1}School of Electrical and Computer Engineering\\
%Georgia Institute of Technology,
%Atlanta, Georgia 30332--0250\\ Email: see http://www.michaelshell.org/contact.html}
%\IEEEauthorblockA{\IEEEauthorrefmark{2}Twentieth Century Fox, Springfield, USA\\
%Email: homer@thesimpsons.com}
%\IEEEauthorblockA{\IEEEauthorrefmark{3}Starfleet Academy, San Francisco, California 96678-2391\\
%Telephone: (800) 555--1212, Fax: (888) 555--1212}
%\IEEEauthorblockA{\IEEEauthorrefmark{4}Tyrell Inc., 123 Replicant Street, Los Angeles, California 90210--4321}}




% use for special paper notices
%\IEEEspecialpapernotice{(Invited Paper)}




% make the title area
\maketitle

% As a general rule, do not put math, special symbols or citations
% in the abstract
\begin{abstract}
Bank transaction fraud is difficult to examine because fraudulent
observations are rare and risk patterns may differ across transaction
contexts. This paper uses exploratory data analysis and visualization to
investigate which transaction characteristics show visible associations
with fraud risk in a synthetic dataset of one million bank transactions.
The analysis combines descriptive summaries, within-group fraud-rate
comparisons, distribution plots, and a merchant-by-hour heatmap. Fraud
accounts for 55,255 transactions, or 5.53\% of the full dataset,
creating a strongly imbalanced outcome. The clearest visual patterns
involve repeated failed attempts, late-night and early-morning timing,
international transaction status, recent PIN changes, and merchant
category. Transactions with two or more failed attempts have
substantially higher fraud rates, while ATM withdrawal, jewelry, and
crypto exchange transactions appear elevated across multiple hours. In
contrast, weekend status and account balance provide weaker or more
cautionary visual separation; location variables were inspected but not
retained as main figures. Visualization therefore identifies candidate
fraud-risk indicators before predictive modeling, although all findings
are associative and based on synthetic data.
\end{abstract}

% keywords
\begin{IEEEkeywords}
bank transaction fraud, exploratory data analysis, data visualization,
class imbalance, fraud rate
\end{IEEEkeywords}

% use for special paper notices



% make the title area
\maketitle

% no keywords

% For peer review papers, you can put extra information on the cover
% page as needed:
% \ifCLASSOPTIONpeerreview
% \begin{center} \bfseries EDICS Category: 3-BBND \end{center}
% \fi
%
% For peerreview papers, this IEEEtran command inserts a page break and
% creates the second title. It will be ignored for other modes.
\IEEEpeerreviewmaketitle


\section{Introduction}\label{introduction}

Transaction fraud monitoring is important because suspicious activity
must be identified within a large volume of legitimate banking
transactions. A central challenge is class imbalance: fraudulent
transactions usually form a small minority, so raw counts can obscure
meaningful differences between transaction groups. A second challenge is
interpretability. Before developing a predictive model, analysts need to
understand whether observable transaction characteristics display stable
and communicable patterns.

Exploratory data analysis (EDA) supports this early stage by summarizing
data quality, revealing skewed distributions, and comparing fraud rates
across behavioural and contextual variables. Visual analysis does not
establish causation, but it can identify candidate indicators for later
statistical testing or predictive modeling. It is also useful for
communication: analysts, risk managers, and other stakeholders can
inspect whether an apparent pattern is consistent across groups or
driven by class size, skewness, or a small subset of observations.

This paper asks: \emph{Which transaction characteristics show visible
associations with fraud risk in bank transaction data?} Its contribution
is a focused visual comparison across authentication behaviour,
transaction timing, international status, recent PIN changes, merchant
context, and selected numerical variables. Because fraud is uncommon,
the analysis emphasizes within-group rates and proportions rather than
raw fraud counts. This paper follows the project proposal's exploratory
scope by prioritizing exploratory visualization over predictive
modeling. Although the proposal mentioned that a simple classifier could
be developed if time permitted, the primary objective was to identify
and explain visible fraud-risk patterns through EDA. The analysis
therefore establishes a reproducible evidence base for describing
candidate variables and interactions without making causal or predictive
claims.

The remainder of the paper reviews relevant fraud-detection and
imbalanced-data research, describes the dataset and reproducible EDA
methodology, reports univariate, bivariate, and multivariate visual
findings, and concludes with practical implications, limitations, and
directions for validation.

\section{Previous Work}\label{previous-work}

Fraud-detection research has examined transaction behaviour,
authentication signals, temporal activity, and contextual variables as
possible inputs to monitoring systems \citep{abdallah2016}.
Anomaly-detection research considers how rare observations can be
distinguished from dominant normal behaviour \citep{chandola2009}.
Research on imbalanced learning also shows why overall accuracy and raw
counts are inadequate when the positive class is uncommon
\citep{he2009}.

The imbalance problem is especially important in financial transaction
data. When legitimate transactions greatly outnumber fraudulent ones, a
method can appear accurate while failing to identify the minority class.
Recent fraud detection research therefore emphasizes resampling,
minority-class synthesis, ensemble methods, and
precision-recall-oriented evaluation
\citep{fiore2019, kennedy2024, bekereci2024}. These studies are
primarily predictive, but their motivation also applies to EDA:
population counts and balanced-sample proportions answer different
questions and must not be interchanged.

Fraud observations may also overlap substantially with legitimate
behaviour. Li \emph{et al.} describe the combination of class imbalance
and class overlap as a central challenge because fraudulent transactions
are often designed to resemble ordinary activity \citep{li2021}. This
supports examining multiple dimensions rather than expecting one
variable to separate the classes. Authentication behaviour, transaction
context, timing, and merchant type may jointly carry more information
than an isolated financial amount.

EDA complements predictive methods by exposing distributions, group
differences, and possible interactions before a model is fitted
\citep{tukey1977}. Bar charts can compare categorical rates,
distribution plots can show overlap, and heatmaps can expose patterns
across contextual dimensions. However, many fraud studies move quickly
to classifier comparison. The present study instead emphasizes a
transparent visual workflow that distinguishes population-rate figures
from balanced-sample figures and reports both strong and weak patterns.
This study remains exploratory: it identifies visible candidate patterns
but does not estimate causal effects or claim predictive performance.

\section{Methodology}\label{methodology}

\subsection{Dataset and Variables}\label{dataset-and-variables}

The full synthetic dataset was obtained from the \emph{Bank Transaction
Fraud Detection Dataset} hosted on Kaggle \citep{kaggleBankFraud}. It
contains 1,000,000 transactions and 26 variables. The binary outcome,
\texttt{is\_fraud}, equals one for a fraudulent transaction and zero
otherwise. The full dataset contains 55,255 fraudulent and 944,745
non-fraudulent transactions, corresponding to an overall fraud rate of
5.53\%.

\begin{table}[!h]
\centering
\caption{\label{tab:table-dataset-summary}Full transaction dataset summary}
\centering
\resizebox{\ifdim\width>\linewidth\linewidth\else\width\fi}{!}{
\begin{tabular}[t]{lrlll}
\toprule
Transactions & Variables & Fraud & Non-fraud & Fraud rate\\
\midrule
1,000,000 & 26 & 55,255 & 944,745 & 5.53\%\\
\bottomrule
\end{tabular}}
\end{table}

Variables examined include transaction type, hour of day, night and
weekend indicators, merchant category, failed attempts, account balance,
distance from home, city, and recent PIN change. Full-data summaries and
fraud-rate calculations are used whenever \texttt{data/transactions.csv}
is available.

The variables can be organized into four conceptual groups. Transaction
context includes domestic or international status, merchant category,
payment method, device type, and location. Temporal context includes
hour, night status, and weekend status. Authentication and account
context includes failed attempts, recent PIN change, account balance,
and credit-related fields. The binary outcome \texttt{is\_fraud}
provides the common basis for comparing these groups.

\subsection{Data Preparation}\label{data-preparation}

Column names were standardized, the binary outcome was converted to
integer form, and readable factors were created for Fraud/Non-fraud,
International/Domestic, Night/Day, Weekend/Weekday, and PIN Changed/No
PIN Change. Data-quality checks covered missingness, repeated
identifiers, duplicate complete rows, and invalid ranges for selected
key variables. In the full dataset, these checks found no missing
values, duplicate transaction IDs, duplicate rows, or invalid values, so
no removal or imputation was required.

A balanced sample containing 1,000 fraud and 1,000 non-fraud
observations is used for selected distributional and sample-composition
figures. Fraud rates by failed attempts and transaction type are
calculated from the full dataset. Because the balanced sample is
deliberately constructed, it cannot estimate population prevalence.

\subsection{Reproducible Analysis
Workflow}\label{reproducible-analysis-workflow}

The paper is written as an executable R Markdown document. The document
reads the CSV data, standardizes variables, calculates every displayed
summary, and generates each figure during rendering. Supporting
functions are stored in a small sourced R file. No static screenshot is
used as an analytical result. This design makes the relationship among
data, calculations, figures, and written interpretations auditable.

The full dataset is used for counts and fraud-rate estimates, including
the failed-attempt comparison in Figure 5. The small balanced sample is
used only for selected distributional or sample-composition figures and
where a runnable submission sample is required. Every sample-based
caption explicitly warns that its shares do not estimate population
prevalence. This separation prevents the 50\% fraud share of the
balanced sample from being confused with the 5.53\% fraud rate in the
full data.

\subsection{Analytical Measures}\label{analytical-measures}

For each category, the descriptive fraud rate is the number of
fraudulent transactions divided by the total number of transactions in
that category. Within-group sample proportions are calculated separately
for fraud and non-fraud cases so that the majority class does not
dominate a distribution plot.

\subsection{Visualization Strategy}\label{visualization-strategy}

Bar charts compare rates across binary and categorical variables. A
faceted point-range plot compares full-data fraud rates across discrete
failed-attempt values and transaction types, with Wilson confidence
intervals used to communicate uncertainty. A heatmap combines hour of
day and merchant category. For grouped full-data comparisons, total
transactions, fraud counts, and fraud rates are calculated directly from
\texttt{df}. Findings are described as visual associations, not causal
effects.

\section{Exploratory Data Analysis
Design}\label{exploratory-data-analysis-design}

The EDA proceeded in three stages. First, univariate checks established
dataset size, class balance, variable ranges, missingness, and duplicate
status. This stage determined whether cleaning or imputation was
necessary and revealed the extent of outcome imbalance. Second,
bivariate comparisons evaluated how fraud rates or within-group
distributions changed across transaction type, time of day, recent PIN
change, weekend status, failed attempts, account balance, and city.
Third, a multivariate heatmap combined merchant category with hour of
day to determine whether two contextual dimensions revealed patterns
hidden by separate marginal summaries.

The visual encodings were selected for the measurement scale and
analytical question. Counts are appropriate for showing class imbalance;
rate bars support fair categorical comparisons; a mosaic plot displays
both transaction-type frequency and within-type composition; faceted
point-and-line displays show fraud rates across discrete failed-attempt
values while preserving transaction context and uncertainty; violin and
box plots show density, median, and overlap; and a heatmap displays a
two-dimensional rate surface. This staged design supports interpretation
without presenting EDA as proof of causality or predictive accuracy.

\section{Results}\label{results}

\subsection{Outcome Imbalance and Dataset
Context}\label{outcome-imbalance-and-dataset-context}

\begin{figure}[!htbp]

{\centering \includegraphics[width=\columnwidth]{Exploratory_Fraud_Risk_Patterns_Vy_Revision_files/figure-latex/fig-class-imbalance-1} 

}

\caption{Class distribution of transaction status. Fraud represents 5.53\% of the full dataset, indicating a highly imbalanced outcome.}\label{fig:fig-class-imbalance}
\end{figure}

Fraud is a minority outcome: the full-data rate is 5.53\%, and the
non-fraud class is more than seventeen times larger. Consequently, later
comparisons use within-group fraud rates rather than raw counts. The
data-quality assessment found no missing values, duplicate transaction
identifiers, duplicate complete rows, or out-of-range values among the
checked fields. Therefore, the observed patterns are not artifacts of an
imputation or row-removal procedure.

\subsection{Transaction Type and Sample
Composition}\label{transaction-type-and-sample-composition}

\begin{figure}[!htbp]

{\centering \includegraphics[width=\columnwidth]{Exploratory_Fraud_Risk_Patterns_Vy_Revision_files/figure-latex/fig-transaction-type-mosaic-1} 

}

\caption{Transaction type and fraud-group membership in the balanced sample. Rectangle widths reflect sample counts; sample shares do not estimate population prevalence.}\label{fig:fig-transaction-type-mosaic}
\end{figure}

Within the balanced sample, international transactions contain a larger
fraud share than domestic transactions. Because fraud and non-fraud
cases were sampled equally, this figure describes sample composition and
is not a population fraud rate. Specifically, fraud represents 68.1\% of
sampled international transactions but 45.1\% of sampled domestic
transactions. The difference is directionally consistent with the
full-data international fraud-rate result, but the percentages in Fig. 2
must remain sample-composition measures.

\subsection{Binary Transaction
Characteristics}\label{binary-transaction-characteristics}

\begin{figure}[!htbp]

{\centering \includegraphics[width=\columnwidth]{Exploratory_Fraud_Risk_Patterns_Vy_Revision_files/figure-latex/fig-binary-characteristics-1} 

}

\caption{Fraud rates across selected binary transaction characteristics. Night transactions and recent PIN changes show higher fraud rates, while weekend status shows little separation.}\label{fig:fig-binary-characteristics}
\end{figure}

Night transactions and transactions following a recent PIN change show
visibly higher fraud rates than their comparison groups. Night
transactions have a 7.96\% fraud rate compared with 4.07\% during the
day. Transactions following a recent PIN change reach 8.36\%, compared
with 5.28\% when no recent change is recorded. Weekend and weekday rates
are nearly equal at 5.57\% and 5.51\%, respectively. This contrast is
useful because it shows that not every intuitive context variable
separates fraud. These differences are candidate visual indicators, not
evidence of causation.

\subsection{Account Balance}\label{account-balance}

The fraud group has a somewhat higher median balance in the balanced
sample, but the distributions overlap substantially. Account balance
therefore provides weaker standalone separation than failed attempts or
transaction context. The sample median is approximately \$8,157 for
fraud and \$7,390 for non-fraud, but the interquartile ranges and violin
densities occupy much of the same region. This is a useful negative
result: a shift in central tendency does not imply that the variable is
an effective standalone discriminator.

\begin{figure}[!htbp]

{\centering \includegraphics[width=0.95\columnwidth]{Exploratory_Fraud_Risk_Patterns_Vy_Revision_files/figure-latex/fig-account-balance-1} 

}

\caption{Account-balance distribution shapes by fraud group in the balanced sample. Violin width shows density and inset boxes show the median and interquartile range; the balanced sample does not estimate population prevalence.}\label{fig:fig-account-balance}
\end{figure}

\subsection{Failed Attempts Across Transaction
Contexts}\label{failed-attempts-across-transaction-contexts}

\begin{figure*}[!t]

{\centering \includegraphics[width=0.94\textwidth]{Exploratory_Fraud_Risk_Patterns_Vy_Revision_files/figure-latex/fig-failed-attempts-1} 

}

\caption{Full-data fraud rates by failed attempts and transaction type. Points show descriptive fraud rates, and vertical bars show 95\% Wilson confidence intervals. The largest increase occurs between one and two failed attempts in both transaction contexts.}\label{fig:fig-failed-attempts}
\end{figure*}

The full-data results show a clear increase at two failed attempts.
Domestic fraud rates rise from about 4\% at zero or one attempt to
roughly 14\% at two or more attempts, while international rates increase
from about 9\% to around 19--20\%. Because both curves follow a similar
shape, the figure suggests a consistent failed-attempt pattern together
with a higher international fraud baseline, rather than a strong
interaction.

\subsection{Merchant and Time-Based
Patterns}\label{merchant-and-time-based-patterns}

\begin{figure*}[!t]

{\centering \includegraphics[width=0.94\textwidth]{Exploratory_Fraud_Risk_Patterns_Vy_Revision_files/figure-latex/fig-multivariate-patterns-1} 

}

\caption{Multivariate fraud patterns: (A) hourly fraud rates for domestic and international transactions and (B) merchant-category fraud rates across hours.}\label{fig:fig-multivariate-patterns}
\end{figure*}

Panel A shows that international transactions have higher fraud rates
than domestic transactions across most hours within this synthetic
dataset. Both series also rise during late-night and early-morning
periods. Panel B adds merchant context: ATM withdrawal, jewelry, and
crypto exchange remain visually elevated across several hours, whereas
most other categories cluster at lower rates. Aggregated across hours,
the three highlighted categories have fraud rates of approximately
8.65--8.74\%, compared with roughly 4.62--4.84\% for most other merchant
categories. Together, the panels show an exploratory pattern in which
transaction timing, international status, and merchant context are
visually associated with different fraud rates. These associations are
candidate indicators within the synthetic dataset, not causal evidence.

\subsection{Summary of Visual
Findings}\label{summary-of-visual-findings}

The strongest visible associations involve repeated failed attempts,
late-night and early-morning timing, international status, recent PIN
changes, and merchant category. Account balance, weekend status, payment
method, device type, and distance from home show weaker separation or
require cautious interpretation.

Additional variables mentioned in the proposal, including payment
method, device type, transaction amount, distance from home, and other
location-related fields, were inspected during EDA; however, they showed
weaker or less central visual separation and are therefore summarized
rather than displayed as main figures within the IEEE page limit.

\FloatBarrier

\section{Discussion}\label{discussion}

\subsection{Interpretation and Practical
Relevance}\label{interpretation-and-practical-relevance}

EDA identifies candidate features, reveals the effect of class
imbalance, and shows which variables have limited standalone separation.
Stronger patterns could inform feature selection, monitoring rules,
interactions, or later predictive experiments. Visual prominence alone
does not establish statistical significance or operational usefulness.

The results suggest a hierarchy of candidate information. Authentication
behaviour is represented by failed attempts and recent PIN change;
transaction context is represented by international status and merchant
category; and temporal context is represented by night or early-morning
activity. These signals are interpretable and can be reviewed by a human
analyst. In contrast, weekend status and account balance provide weaker
separation. Including weak results is important because an EDA should
narrow the candidate set rather than merely collect visually interesting
charts.

For an operational fraud-monitoring workflow, the findings could
motivate rules or model features such as an indicator for two or more
failed attempts, an interaction between merchant category and hour, or
additional review when a recent PIN change coincides with an unusual
context. Such uses would require validation on out-of-sample real
transactions, threshold selection based on false-positive costs, and
continuous monitoring for changing fraud behaviour. The figures should
therefore be viewed as hypothesis-generating tools, not as decision
rules.

\subsection{Limitations}\label{limitations}

The dataset is synthetic, so the findings describe patterns within this
dataset rather than verified real-world banking behaviour. Location
variables are treated as categorical features, not validated geographic
relationships. EDA cannot establish causation, and deliberately balanced
sample plots cannot estimate population prevalence. The associations
should be validated with statistical analysis and predictive evaluation
that accounts for class imbalance.

The analysis is also cross-sectional. It does not examine how
relationships change across days, customers, or evolving fraud
strategies. Aggregated rates may hide repeated activity by the same
account. Finally, no classifier, train-test split, or cost-sensitive
evaluation is included because the proposal identified modeling only as
an optional extension if time permitted, while the primary deliverable
remained EDA and visualization.

\section{Conclusions}\label{conclusions}

This paper used code-generated exploratory visualizations to examine
which bank transaction characteristics show visible associations with
fraud risk. Fraud represents 5.53\% of the full dataset, making
rate-based comparisons essential. Repeated failed attempts, night and
early-morning timing, international status, recent PIN changes, and
merchant category provide the clearest visual patterns. Weekend status
and account balance are weaker or more cautionary findings.

The study demonstrates why visual design and sampling language matter in
an imbalanced problem. Full-data rates describe prevalence, while
balanced-sample figures describe group composition and distributional
differences. Maintaining this distinction prevents overstated
conclusions and makes the analysis reproducible.

Future work should validate these exploratory patterns using independent
or real transaction data and examine whether the visible relationships
remain stable over time. Predictive modeling may be considered as a
separate extension after the exploratory findings are validated; it is
not part of the present paper's contribution.

\renewcommand\refname{References}
\bibliography{support/references.bib}
\end{document}

